\appendix
\section{小组分工}

\begin{table}[H]
    \centering
    \caption{实验小组成员分工}
    \begin{tabular}{@{}llp{7cm}@{}}
        \toprule
        \textbf{姓名} & \textbf{学号} & \textbf{负责内容} \\
        \midrule
        江李阳 & 202400460042 & 协议分析设计,代码撰写,实验报告撰写 \\
        王健一 & 202400460037 & 审计实验报告,修改整合与润色实验报告 \\
        邱泽辉 & 202400460046 & 协议分析设计与调研,代码撰写\\
        \bottomrule
    \end{tabular}
\end{table}

\section{AI 使用声明}

本报告在编写过程中使用了人工智能辅助工具，具体使用情况声明如下：

\subsection{使用工具信息}
\begin{itemize}[leftmargin=2em]
    \item \textbf{工具名称}:claude code cli
    \item \textbf{模型版本}: deepseek v4 pro
    \item \textbf{用途范围}:
    \begin{itemize}[leftmargin=2em]
        \item 实验报告 \LaTeX{} 排版调试、语法纠错与格式优化
        \item 代码片段的提取、整理与嵌入文档（\texttt{lstlisting} 格式化）
        \item 项目代码的调试修改
    \end{itemize}
\end{itemize}

\subsection{生成内容说明}
\begin{enumerate}[leftmargin=2em]
    \item \textbf{文本内容}:实验原理描述、协议步骤分析、结果分析章节由作者独立撰写，AI 辅助进行语句润色与格式调整。
    \item \textbf{图表数据}:所有测试截图、性能数据均来自真实运行环境的人工操作记录，AI 未参与任何数据生成或伪造。
\end{enumerate}

% \subsection{责任申明}
% 本人对报告全部内容（含经 AI 辅助润色或格式化的部分）的真实性、准确性与原创性负全部责任。本报告中不存在由 AI 生成的虚假引用、编造数据或伪造结果。如有违反学术诚信的行为，愿承担相应后果。

% \vspace{1em}
% \begin{center}
% \rule{0.5\textwidth}{0.4pt}\\[0.5em]
% 签名：\underline{\hspace{5cm}} \quad 日期：\underline{\hspace{3cm}}
% \end{center}

\section{项目完整代码地址}
\url{https://github.com/hxhdhcjmet/a2aProtocal.git}